% !PRINTED 222
\weberpage{222}
%<<CONT
that $b = \chi(a)$, $b' = \chi(a')$ are the two poles of the substitution
$\chi^{-1}\Theta\chi$. % [sic/mathematics: with b=chi(a), theorem 53.2 requires chi Theta chi^{-1}; the source prints the inverse conjugation order and the ensuing calculation follows that printed order.]
In this we can take $\infty$ and 0 for $a$ and $a'$. and $\omega$ for $b$. % [sic: the source prints a full stop before the lower-case ``und''; preserved here rather than silently regularized.]
Then $b'$ becomes a pole $\xi$ from the sequence (7),
still to be determined more closely, and for $\chi$ we obtain the form

\[
\chi(x) = \frac{\alpha x + \beta}{\gamma x + \delta} = \begin{pmatrix} \alpha, & \beta \\ \gamma, & \delta \end{pmatrix},
\]

in which

\weq{8}{\omega = \frac{\alpha}{\gamma}, \quad \xi = \frac{\beta}{\delta}}

is to be set.
Now we can take for $\Theta$ any of the substitutions $\varepsilon^h x$,
provided only that $h$ is not divisible by 5, and thus obtain the substitution


\[
\begin{aligned}
\chi^{-1}\Theta\chi &= \begin{pmatrix} \delta, & -\beta \\ -\gamma, & \alpha \end{pmatrix}\begin{pmatrix} \varepsilon^h, & 0 \\ 0, & 1 \end{pmatrix}\begin{pmatrix} \alpha, & \beta \\ \gamma, & \delta \end{pmatrix} \\
&= \begin{pmatrix} \varepsilon^h\alpha\delta - \beta\gamma, & (\varepsilon^h-1)\beta\delta \\ -(\varepsilon^h-1)\alpha\gamma, & \alpha\delta - \varepsilon^h\beta\gamma \end{pmatrix},
\end{aligned}
\]

whose poles must be $\omega$ and $\xi$. By \S.\,51, (6), then, $\omega$ and $\xi$
must be the roots of the quadratic equation

\[
\alpha\gamma x^2 + (\alpha\delta + \beta\gamma)x + \beta\delta = 0
\]

from which

\[
\frac{\beta\delta}{\alpha\gamma} = \omega\xi, \quad \frac{\delta}{\gamma} + \frac{\beta}{\alpha} = -\omega - \xi
\]

follows. From the second of these equations there results by means of (8):

\[
(\alpha + \delta)(\alpha\delta + \beta\gamma) = 0,
\]

and since $\alpha\delta + \beta\gamma$ cannot vanish, because otherwise $\omega + \xi = 0$
would have to hold, which by (7) is not possible, we have $\alpha + \delta = 0$.
Accordingly we obtain for $\chi$, if for the sake of simplicity we assume
$\gamma = 1$:

\weq{9}{\chi = \begin{pmatrix} \omega, & \beta \\ 1, & -\omega \end{pmatrix}.}

Here $\beta$, which by (8) has the value $-\omega\xi$, is still to be
determined.

If we assume the substitution $\chi$ as known, then
we can form the whole icosahedral group.

Namely, if one lets $r$ and $s$ run through the numbers 0, 1, 2, 3, 4,
one has in this group the substitutions

\weq{10}{\Theta^r, \quad \psi\Theta^r, \quad \Theta^r\chi\Theta^s, \quad \Theta^r\chi\psi\Theta^s,}
%CONT>>
